\documentclass[11pt]{article}
\usepackage[margin=1in]{geometry}
\usepackage{amsmath, amssymb, amsthm}
\usepackage{enumitem}
\usepackage{xcolor}
\usepackage{tcolorbox}
\tcbuselibrary{skins, breakable}

% Define color boxes for different framework sections
\definecolor{problemconfig}{RGB}{230, 242, 255}
\definecolor{strategic}{RGB}{255, 230, 242}
\definecolor{tactical}{RGB}{242, 255, 230}
\definecolor{techniques}{RGB}{255, 242, 230}
\definecolor{verification}{RGB}{255, 230, 230}
\definecolor{solution}{RGB}{230, 255, 255}
\definecolor{competition}{RGB}{242, 242, 255}
\definecolor{gold}{RGB}{218, 165, 32}

\newtcolorbox{problemconfigbox}{
    colback=problemconfig, colframe=blue!40!black, title=PROBLEM CONFIGURATION ANALYSIS,
    fonttitle=\bfseries, breakable
}

\newtcolorbox{strategicbox}{
    colback=strategic, colframe=pink!40!black, title=STRATEGIC FRAMEWORK APPLICATION,
    fonttitle=\bfseries, breakable
}

\newtcolorbox{tacticalbox}{
    colback=tactical, colframe=green!40!black, title=TACTICAL METHODOLOGY SELECTION,
    fonttitle=\bfseries, breakable
}

\newtcolorbox{techniquesbox}{
    colback=techniques, colframe=orange!40!black, title=CONVERSION TECHNIQUES APPLICATION,
    fonttitle=\bfseries, breakable
}

\newtcolorbox{verificationbox}{
    colback=verification, colframe=red!40!black, title=VERIFICATION STRATEGY DESIGN,
    fonttitle=\bfseries, breakable
}

\newtcolorbox{solutionbox}{
    colback=solution, colframe=cyan!40!black, title=SOLUTION ROADMAP,
    fonttitle=\bfseries, breakable
}

\newtcolorbox{competitionbox}{
    colback=competition, colframe=purple!40!black, title=COMPETITION STRATEGY INTEGRATION,
    fonttitle=\bfseries, breakable
}

\newtcolorbox{keyinsightbox}{
    colback=yellow!20, colframe=gold, title=KEY INSIGHT,
    fonttitle=\bfseries, arc=3pt, boxrule=2pt, leftrule=5pt
}

\newtcolorbox{warningbox}{
    colback=red!10, colframe=red!50, title=WARNING: Common Pitfall,
    fonttitle=\bfseries, arc=3pt, boxrule=2pt, leftrule=5pt
}

\newtcolorbox{tipbox}{
    colback=green!10, colframe=green!50, title=PRO TIP,
    fonttitle=\bfseries, arc=3pt, boxrule=2pt, leftrule=5pt
}

\title{\textbf{2017 AMC 12B Problem 20: Strategic Analysis}}
\author{Using AMC12/COMC Problem-Solving Framework}
\date{}

\begin{document}
\maketitle

\section*{Problem Statement}
Real numbers $x$ and $y$ are chosen independently and uniformly at random from the interval $(0,1)$. What is the probability that $\lfloor\log_2x\rfloor=\lfloor\log_2y\rfloor$?

\textbf{Answer Choices:}\\
$\textbf{(A)}\ \frac{1}{8}\qquad\textbf{(B)}\ \frac{1}{6}\qquad\textbf{(C)}\ \frac{1}{4}\qquad\textbf{(D)}\ \frac{1}{3}\qquad\textbf{(E)}\ \frac{1}{2}$

\begin{problemconfigbox}
\subsection*{A. Problem Classification}
\begin{itemize}[leftmargin=*]
    \item \textbf{Problem Type}: To Find -- determine a probability value
    \item \textbf{Difficulty Band}: Late-band (Problem 20 of 25) -- requires understanding of logarithms, floor functions, and geometric series
    \item \textbf{Competition Context}: AMC12 (75 min, 25 problems) -- approximately 3-4 minutes allocated for late-band problems
    \item \textbf{Mathematical Domain}: Probability with strong connections to logarithms, floor functions, and infinite series
\end{itemize}

\subsection*{B. Problem Structure Identification}
\begin{itemize}[leftmargin=*]
    \item \textbf{Given Information}: 
    \begin{itemize}
        \item $x, y$ chosen independently and uniformly from $(0,1)$
        \item Need probability that $\lfloor\log_2 x\rfloor = \lfloor\log_2 y\rfloor$
    \end{itemize}
    
    \item \textbf{Constraints}: 
    \begin{itemize}
        \item Open interval $(0,1)$ means $0 < x, y < 1$
        \item For $0 < x < 1$: $\log_2 x < 0$ (always negative)
        \item Floor function creates discrete integer values from continuous logarithms
        \item Independence: joint probability is product of marginal probabilities
    \end{itemize}
    
    \item \textbf{Objective}: Find $P(\lfloor\log_2 x\rfloor = \lfloor\log_2 y\rfloor)$
    
    \item \textbf{Hidden Assumptions}: 
    \begin{itemize}
        \item The interval $(0,1)$ partitions naturally based on floor values
        \item Each floor value corresponds to a dyadic interval: $[\frac{1}{2^{n+1}}, \frac{1}{2^n})$
        \item These intervals have lengths that form a geometric sequence
        \item The probability involves an infinite series
    \end{itemize}
    
    \item \textbf{Answer Choice Leverage}: 
    \begin{itemize}
        \item All choices are unit fractions with small denominators
        \item $\frac{1}{2}$ would be too high (requires strong correlation)
        \item $\frac{1}{8}$ seems too low (there's overlap)
        \item $\frac{1}{3}$ and $\frac{1}{4}$ are geometric series sums ($\frac{1}{1-r}$ pattern)
        \item Can eliminate extreme values and test middle options
    \end{itemize}
\end{itemize}

\subsection*{C. Strategic Orientation}
\begin{itemize}[leftmargin=*]
    \item \textbf{Hypothesis}: Let $X, Y \sim \text{Uniform}(0,1)$ be independent
    
    \item \textbf{Conclusion}: Determine $P(\lfloor\log_2 X\rfloor = \lfloor\log_2 Y\rfloor)$
    
    \item \textbf{Notation}: 
    \begin{itemize}
        \item $X, Y$: random variables uniformly distributed on $(0,1)$
        \item $\lfloor \cdot \rfloor$: floor function (greatest integer function)
        \item $\log_2$: logarithm base 2
        \item For $x \in (0,1)$: $\lfloor\log_2 x\rfloor \in \{-1, -2, -3, \ldots\}$
    \end{itemize}
    
    \item \textbf{Intuitive Approach}: 
    \begin{itemize}
        \item Partition $(0,1)$ by floor values of $\log_2$
        \item $\lfloor\log_2 x\rfloor = -1$ when $x \in [\frac{1}{2}, 1)$ (length $\frac{1}{2}$)
        \item $\lfloor\log_2 x\rfloor = -2$ when $x \in [\frac{1}{4}, \frac{1}{2})$ (length $\frac{1}{4}$)
        \item Pattern: each interval has length $\frac{1}{2^n}$
        \item Sum probabilities that both fall in same interval
    \end{itemize}
    
    \item \textbf{Cross-Domain Potential}: 
    \begin{itemize}
        \item \textbf{Geometric}: Visualize as areas in coordinate plane
        \item \textbf{Algebraic}: Use series summation formulas
        \item \textbf{Analysis}: Understand interval structure through logarithm properties
    \end{itemize}
\end{itemize}
\end{problemconfigbox}

\begin{strategicbox}
\subsection*{A. Psychological Strategies}
\begin{itemize}[leftmargin=*]
    \item \textbf{Mental Toughness}: Don't be intimidated by the nested functions (floor of logarithm); break down step-by-step
    
    \item \textbf{Creativity Cultivation}: Recognize that floor function creates discrete regions from continuous domain; visualize interval partitioning
    
    \item \textbf{Iterative Refinement}: Start with small cases ($n=-1, -2$), identify pattern, then generalize to infinite series
\end{itemize}

\subsection*{B. Investigation Strategies}
\begin{itemize}[leftmargin=*]
    \item \textbf{Startup Orientation}: Recognize compound function structure: probability $\rightarrow$ floor function $\rightarrow$ logarithm
    
    \item \textbf{Get Your Hands Dirty}: Test specific values to understand floor behavior
    \begin{itemize}
        \item $x = 0.75 \in [\frac{1}{2}, 1)$: $\log_2(0.75) \approx -0.415$, so $\lfloor\log_2(0.75)\rfloor = -1$
        \item $x = 0.3 \in [\frac{1}{4}, \frac{1}{2})$: $\log_2(0.3) \approx -1.737$, so $\lfloor\log_2(0.3)\rfloor = -2$
    \end{itemize}
    
    \item \textbf{Penultimate Step}: Once we identify the geometric series pattern, the final step is applying the formula $\sum_{n=0}^{\infty} ar^n = \frac{a}{1-r}$
    
    \item \textbf{Wishful Thinking}: We wish for a simple pattern in interval lengths -- and indeed, powers of $\frac{1}{2}$ emerge!
    
    \item \textbf{Make It Easier}: Instead of thinking about logarithms, think about powers of 2: $2^{-1} = \frac{1}{2}$, $2^{-2} = \frac{1}{4}$, etc.
    
    \item \textbf{Conjecture via Small Cases}: 
    \begin{itemize}
        \item Case $n=-1$: $P(\text{both in } [\frac{1}{2},1)) = (\frac{1}{2})^2 = \frac{1}{4}$
        \item Case $n=-2$: $P(\text{both in } [\frac{1}{4},\frac{1}{2})) = (\frac{1}{4})^2 = \frac{1}{16}$
        \item Pattern: probabilities are $\frac{1}{4}, \frac{1}{16}, \frac{1}{64}, \ldots$ (geometric with ratio $\frac{1}{4}$)
    \end{itemize}
    
    \item \textbf{Visualization}: Draw number line $(0,1)$ partitioned at $\frac{1}{2}, \frac{1}{4}, \frac{1}{8}, \ldots$; or draw unit square $(0,1) \times (0,1)$ with diagonal stripes for matching floor values
\end{itemize}

\subsection*{C. Late-Band Strategic Playbook}
\begin{itemize}[leftmargin=*]
    \item \textbf{Answer-Choice Leverage}: Recognize $\frac{1}{3}$ suggests geometric series with $r = \frac{1}{4}$ (since $\frac{a}{1-r} = \frac{1/4}{1-1/4} = \frac{1/4}{3/4} = \frac{1}{3}$)
    
    \item \textbf{Make It Easier}: Convert logarithm problem to interval problem; convert continuous to discrete (via floor)
    
    \item \textbf{Penultimate Step}: Identify that answer is $\sum_{n=1}^{\infty} (\frac{1}{2^n})^2 = \sum_{n=1}^{\infty} \frac{1}{4^n}$
    
    \item \textbf{Wishful Thinking}: Hope for geometric series -- and it appears!
    
    \item \textbf{Conjecture via Small Cases}: Compute first few terms, recognize pattern, apply formula
\end{itemize}

\subsection*{D. Argument Construction Strategy}
\begin{itemize}[leftmargin=*]
    \item \textbf{Direct Proof}: 
    \begin{enumerate}
        \item Partition $(0,1)$ by floor values
        \item Compute probability for each case
        \item Sum using geometric series formula
    \end{enumerate}
    
    \item \textbf{Proof by Contradiction}: Not applicable
    
    \item \textbf{Mathematical Induction}: Could verify pattern, but direct summation is cleaner
    
    \item \textbf{Stylistic Conventions}: Use case analysis; leverage independence and uniform distribution
\end{itemize}

\subsection*{E. Cross-Domain Translation}
\begin{itemize}[leftmargin=*]
    \item \textbf{Draw a Picture}: Unit square with diagonal stripes (regions where floor values match)
    
    \item \textbf{Recast the Problem}: ``Interval matching'' problem rather than ``floor of logarithm'' problem
    
    \item \textbf{Change Your Point of View}: Think in terms of powers of 2 rather than logarithms
\end{itemize}
\end{strategicbox}

\begin{tacticalbox}
\subsection*{A. Core Mathematical Tactics}
\begin{itemize}[leftmargin=*]
    \item \textbf{Symmetry}: Independence and identical distributions create symmetry
    
    \item \textbf{Extreme Principle}: Not directly applicable
    
    \item \textbf{Invariants}: Floor value partitions are invariant under the problem structure
    
    \item \textbf{Monovariants}: Not applicable
    
    \item \textbf{Parity}: Not applicable
    
    \item \textbf{Graph Theory}: Not applicable
    
    \item \textbf{Generating Functions}: Not primary, though series summation is related
    
    \item \textbf{Pattern Recognition}: Identify geometric sequence in interval lengths and probabilities
\end{itemize}

\subsection*{B. Crossover Techniques}
\begin{itemize}[leftmargin=*]
    \item \textbf{Algebraic-Probabilistic}: Convert probability question to series summation
    
    \item \textbf{Geometric-Probabilistic}: Use geometric interpretation via unit square regions
    
    \item \textbf{Analytic-Discrete}: Floor function converts continuous logarithm to discrete integer values
\end{itemize}

\subsection*{C. Domain-Specific Approaches}
\begin{itemize}[leftmargin=*]
    \item \textbf{Probability}: Independence, uniform distribution, total probability (sum over disjoint cases)
    
    \item \textbf{Analysis}: Logarithm properties, floor function behavior, interval decomposition
    
    \item \textbf{Algebra}: Geometric series summation
    
    \item \textbf{Geometry}: Unit square representation (optional but helpful for visualization)
\end{itemize}
\end{tacticalbox}

\begin{techniquesbox}
\subsection*{A. Technique Recognition Index}
\begin{itemize}[leftmargin=*]
    \item \textbf{T17-GeoProb}: Geometric Probability -- applicable for unit square visualization
    \item \textbf{Series Summation}: Geometric series formula -- PRIMARY technique
    \item \textbf{Case Analysis}: Partition by floor values -- essential structure
    \item \textbf{Logarithm Properties}: Understanding $\log_2 x$ for $x \in (0,1)$
\end{itemize}

\subsection*{B. Technique Selection}
\textbf{Primary Technique: Case Analysis + Geometric Series}

\textbf{Why it applies}:
\begin{itemize}
    \item Floor function naturally partitions $(0,1)$ into countably many intervals
    \item Each interval corresponds to one floor value
    \item Interval lengths form geometric sequence: $\frac{1}{2}, \frac{1}{4}, \frac{1}{8}, \ldots$
    \item Probabilities form geometric sequence: $\frac{1}{4}, \frac{1}{16}, \frac{1}{64}, \ldots$
\end{itemize}

\textbf{Configuration cues}: Floor of logarithm with uniform random variables

\textbf{Objective cues}: Finding probability of equality

\textbf{Supporting Technique: Geometric Series Formula}

For $|r| < 1$: $\sum_{n=0}^{\infty} ar^n = \frac{a}{1-r}$

In our case: $a = \frac{1}{4}$, $r = \frac{1}{4}$, giving $\frac{1/4}{1-1/4} = \frac{1/4}{3/4} = \frac{1}{3}$

\subsection*{C. Technique Integration}
\begin{enumerate}
    \item \textbf{Understand domain}: For $x \in (0,1)$, $\log_2 x$ ranges from $-\infty$ to $0$
    \item \textbf{Identify intervals}: $\lfloor\log_2 x\rfloor = n$ when $x \in [2^n, 2^{n+1})$ for $n \in \{\ldots, -3, -2, -1\}$
    \item \textbf{For our interval}: 
    \begin{itemize}
        \item $\lfloor\log_2 x\rfloor = -1$ when $x \in [\frac{1}{2}, 1)$ (length $\frac{1}{2}$)
        \item $\lfloor\log_2 x\rfloor = -2$ when $x \in [\frac{1}{4}, \frac{1}{2})$ (length $\frac{1}{4}$)
        \item $\lfloor\log_2 x\rfloor = -k$ when $x \in [\frac{1}{2^k}, \frac{1}{2^{k-1}})$ (length $\frac{1}{2^k}$)
    \end{itemize}
    \item \textbf{Compute probability}: For both $x$ and $y$ in same interval with length $\ell$: $P = \ell \cdot \ell = \ell^2$
    \item \textbf{Sum over all cases}: $\sum_{k=1}^{\infty} (\frac{1}{2^k})^2 = \sum_{k=1}^{\infty} \frac{1}{4^k} = \frac{1/4}{1-1/4} = \frac{1}{3}$
\end{enumerate}

\subsection*{D. Conflict Resolution}
No conflicts arise. The approach is:
\begin{itemize}[leftmargin=*]
    \item Cases are mutually exclusive (disjoint intervals)
    \item Cases are exhaustive (cover all of $(0,1)$)
    \item Probabilities properly computed via independence
    \item Series converges (geometric with $|r| = \frac{1}{4} < 1$)
\end{itemize}
\end{techniquesbox}

\begin{keyinsightbox}
\textbf{The Deep Insight}:

The key insight is recognizing that the floor function partitions $(0,1)$ into dyadic intervals (powers of 1/2):
$[\frac{1}{2}, 1), [\frac{1}{4}, \frac{1}{2}), [\frac{1}{8}, \frac{1}{4}), \ldots$

Each interval has length $\frac{1}{2^k}$ for $k = 1, 2, 3, \ldots$

The probability that both $x$ and $y$ fall in the same interval is:
$P = \sum_{k=1}^{\infty} \left(\frac{1}{2^k}\right)^2 = \sum_{k=1}^{\infty} \frac{1}{4^k} = \frac{1/4}{1-1/4} = \frac{1}{3}$

This is a geometric series with first term $\frac{1}{4}$ and common ratio $\frac{1}{4}$.

\textbf{Why This is Elegant}:

This problem beautifully connects three mathematical concepts:
\begin{itemize}
    \item \textbf{Logarithms}: Map multiplicative structure (powers of 2) to additive structure (integers)
    \item \textbf{Floor function}: Discretizes continuous range into countable cases
    \item \textbf{Geometric series}: The dyadic structure naturally produces a geometric sequence
\end{itemize}

The base-2 logarithm is particularly natural here because it reveals the binary structure of $(0,1)$, where each ``bit position'' corresponds to one dyadic interval.
\end{keyinsightbox}

\begin{verificationbox}
\subsection*{A. Verification Level Selection}
This is Problem 20 (late-band), so use \textbf{Level 3: Comprehensive Verification}

\textbf{Decision criteria}:
\begin{itemize}
    \item High difficulty (P20/25)
    \item Multiple conceptual steps (logarithms, floor, probability, series)
    \item Requires correct understanding of interval structure
    \item Series summation must be accurate
\end{itemize}

\subsection*{B. Specialized Verification Methods}
\begin{itemize}[leftmargin=*]
    \item \textbf{Interval Check}:
    \begin{itemize}
        \item Verify $\lfloor\log_2(\frac{1}{2})\rfloor = \lfloor -1 \rfloor = -1$ \checkmark
        \item Verify $\lfloor\log_2(\frac{1}{4})\rfloor = \lfloor -2 \rfloor = -2$ \checkmark
        \item Verify intervals partition $(0,1)$: $\sum_{k=1}^{\infty} \frac{1}{2^k} = 1$ \checkmark
    \end{itemize}
    
    \item \textbf{Series Computation}:
    \begin{itemize}
        \item First term: $(\frac{1}{2})^2 = \frac{1}{4}$ \checkmark
        \item Common ratio: $\frac{1}{16} \div \frac{1}{4} = \frac{1}{4}$ \checkmark
        \item Sum: $\frac{1/4}{1-1/4} = \frac{1/4}{3/4} = \frac{1}{3}$ \checkmark
    \end{itemize}
    
    \item \textbf{Boundary Cases}:
    \begin{itemize}
        \item What if both in $[\frac{1}{2}, 1)$? Probability $= \frac{1}{4}$ (first term) \checkmark
        \item As $k \to \infty$, terms $\to 0$ (series converges) \checkmark
    \end{itemize}
    
    \item \textbf{Answer Choice Verification}:
    \begin{itemize}
        \item Check answer $\frac{1}{3}$: This equals geometric series with $a = \frac{1}{4}$, $r = \frac{1}{4}$ \checkmark
        \item Sanity check: $0 < \frac{1}{3} < \frac{1}{2}$ (probability is positive and less than $\frac{1}{2}$) \checkmark
    \end{itemize}
    
    \item \textbf{Alternative Verification via Geometric Probability}:
    \begin{itemize}
        \item Draw unit square $(0,1) \times (0,1)$
        \item Favorable region: union of squares $[\frac{1}{2^{k+1}}, \frac{1}{2^k}) \times [\frac{1}{2^{k+1}}, \frac{1}{2^k})$
        \item Total area: $\sum_{k=1}^{\infty} (\frac{1}{2^k})^2 = \frac{1}{3}$ \checkmark
    \end{itemize}
\end{itemize}

\subsection*{C. Confidence Calibration}
\textbf{Confidence Level}: 95\% (High)

\textbf{Reasoning}:
\begin{itemize}
    \item Interval structure is correct (verified with specific values)
    \item Series computation is standard and reliable
    \item Answer matches expected form (unit fraction from geometric series)
    \item Multiple verification methods agree
\end{itemize}

\textbf{Potential concerns}: None significant; approach is sound and well-verified
\end{verificationbox}

\begin{warningbox}
\textbf{Common Pitfall \#1}: Forgetting that $\log_2 x < 0$ for $x \in (0,1)$. Students may incorrectly think floor values are positive.

\textbf{Common Pitfall \#2}: Misidentifying interval boundaries. Remember:
\begin{itemize}
    \item $\lfloor\log_2 x\rfloor = -1$ for $x \in [\frac{1}{2}, 1)$, NOT $[\frac{1}{2}, 1]$
    \item Interval is half-open: $[2^{-k}, 2^{-k+1})$
\end{itemize}

\textbf{Common Pitfall \#3}: Incorrect geometric series formula. Be careful with indexing:
\begin{itemize}
    \item $\sum_{n=0}^{\infty} ar^n = \frac{a}{1-r}$ (starting from $n=0$)
    \item $\sum_{n=1}^{\infty} ar^n = \frac{ar}{1-r}$ (starting from $n=1$)
\end{itemize}

In our problem: $\sum_{k=1}^{\infty} \frac{1}{4^k}$ can be computed as:
\begin{itemize}
    \item Method 1: $\sum_{k=1}^{\infty} 1 \cdot (\frac{1}{4})^k = \frac{1 \cdot \frac{1}{4}}{1-\frac{1}{4}} = \frac{1/4}{3/4} = \frac{1}{3}$ $\checkmark$
    \item Method 2: $\sum_{k=0}^{\infty} (\frac{1}{4})^k - (\frac{1}{4})^0 = \frac{1}{1-\frac{1}{4}} - 1 = \frac{4}{3} - 1 = \frac{1}{3}$ $\checkmark$
\end{itemize}

\textbf{Common Pitfall \#4}: Thinking the problem asks for $P(x = y)$ (probability of exact equality). The floor function makes this a discrete matching problem, not a continuous one.
\end{warningbox}

\begin{tipbox}
\textbf{Pro Tip \#1}: When you see floor of logarithm, immediately think about interval partitioning by powers of the base.

\textbf{Pro Tip \#2}: For $x \in (0,1)$, $\log_b x$ is always negative. The floor values are $-1, -2, -3, \ldots$

\textbf{Pro Tip \#3}: Dyadic intervals (powers of 2) are natural for base-2 logarithms. This creates geometric sequences in both interval lengths and probabilities.

\textbf{Pro Tip \#4}: Geometric series with ratio $r = \frac{1}{4}$ gives sum $\frac{a}{3/4} = \frac{4a}{3}$. If first term is $\frac{1}{4}$, sum is $\frac{1}{3}$.

\textbf{Pro Tip \#5}: Visualize this problem as a unit square $[0,1] \times [0,1]$ in the $(x,y)$ plane. The favorable region consists of small squares along the ``diagonal'' where floor values match:
\begin{itemize}
    \item Square 1: $[\frac{1}{2}, 1) \times [\frac{1}{2}, 1)$ with area $\frac{1}{4}$
    \item Square 2: $[\frac{1}{4}, \frac{1}{2}) \times [\frac{1}{4}, \frac{1}{2})$ with area $\frac{1}{16}$
    \item Square 3: $[\frac{1}{8}, \frac{1}{4}) \times [\frac{1}{8}, \frac{1}{4})$ with area $\frac{1}{64}$
    \item And so on...
\end{itemize}
Total favorable area = $\frac{1}{4} + \frac{1}{16} + \frac{1}{64} + \cdots = \frac{1}{3}$
\end{tipbox}

\begin{solutionbox}
\subsection*{Solution Roadmap}

\textbf{Step 1: Understand the Floor Function Behavior}
\begin{itemize}
    \item For $x \in (0,1)$: $\log_2 x < 0$ (negative)
    \item $\lfloor\log_2 x\rfloor$ takes values $\ldots, -3, -2, -1$
    \item $\lfloor\log_2 x\rfloor = n \iff n \leq \log_2 x < n+1 \iff 2^n \leq x < 2^{n+1}$
\end{itemize}

\textbf{Step 2: Partition $(0,1)$ by Floor Values}

\textit{Key Understanding}: $\lfloor\log_2 x\rfloor = n$ means $n \leq \log_2 x < n+1$, which is equivalent to $2^n \leq x < 2^{n+1}$

\begin{itemize}
    \item $\lfloor\log_2 x\rfloor = -1$ when $x \in [\frac{1}{2}, 1)$ 
    \begin{itemize}
        \item Derivation: $-1 \leq \log_2 x < 0 \iff 2^{-1} \leq x < 2^0 \iff \frac{1}{2} \leq x < 1$
        \item Check: $\log_2(1/2) = -1$, $\log_2(1) = 0$ $\checkmark$
        \item Interval length: $1 - \frac{1}{2} = \frac{1}{2}$
        \item Probability: $P(x \in [\frac{1}{2}, 1)) = \frac{1}{2}$
    \end{itemize}
    
    \item $\lfloor\log_2 x\rfloor = -2$ when $x \in [\frac{1}{4}, \frac{1}{2})$
    \begin{itemize}
        \item Derivation: $-2 \leq \log_2 x < -1 \iff 2^{-2} \leq x < 2^{-1} \iff \frac{1}{4} \leq x < \frac{1}{2}$
        \item Check: $\log_2(1/4) = -2$, $\log_2(1/2) = -1$ $\checkmark$
        \item Interval length: $\frac{1}{2} - \frac{1}{4} = \frac{1}{4}$
        \item Probability: $P(x \in [\frac{1}{4}, \frac{1}{2})) = \frac{1}{4}$
    \end{itemize}
    
    \item \textbf{General Pattern}: $\lfloor\log_2 x\rfloor = -k$ when $x \in [\frac{1}{2^k}, \frac{1}{2^{k-1}})$ 
    \begin{itemize}
        \item Interval length: $\frac{1}{2^{k-1}} - \frac{1}{2^k} = \frac{2 - 1}{2^k} = \frac{1}{2^k}$
        \item Probability: $P(x \in [\frac{1}{2^k}, \frac{1}{2^{k-1}})) = \frac{1}{2^k}$
    \end{itemize}
    
    \item \textbf{Coverage Check}: These intervals partition $(0,1)$ completely:
    \begin{align*}
    \sum_{k=1}^{\infty} \frac{1}{2^k} &= \frac{1/2}{1-1/2} = \frac{1/2}{1/2} = 1 \checkmark
    \end{align*}
\end{itemize}

\textbf{Step 3: Compute Probability for Each Case}
\begin{itemize}
    \item Since $x$ and $y$ are independent and uniform on $(0,1)$:
    \item $P(\lfloor\log_2 x\rfloor = \lfloor\log_2 y\rfloor = -k)$
    \item $= P(x \in [\frac{1}{2^k}, \frac{1}{2^{k-1}})) \cdot P(y \in [\frac{1}{2^k}, \frac{1}{2^{k-1}}))$
    \item $= \frac{1}{2^k} \cdot \frac{1}{2^k} = \frac{1}{4^k}$
\end{itemize}

\textbf{Specific Cases:}
\begin{itemize}
    \item $k=1$: $P(\text{both in } [\frac{1}{2}, 1)) = (\frac{1}{2})^2 = \frac{1}{4}$
    \item $k=2$: $P(\text{both in } [\frac{1}{4}, \frac{1}{2})) = (\frac{1}{4})^2 = \frac{1}{16}$
    \item $k=3$: $P(\text{both in } [\frac{1}{8}, \frac{1}{4})) = (\frac{1}{8})^2 = \frac{1}{64}$
\end{itemize}

\textbf{Step 4: Sum Over All Cases}
\begin{itemize}
    \item By law of total probability (cases are disjoint):
    \begin{align*}
    P(\lfloor\log_2 x\rfloor = \lfloor\log_2 y\rfloor) &= \sum_{k=1}^{\infty} P(\text{both have floor value } -k)\\
    &= \sum_{k=1}^{\infty} \frac{1}{4^k}\\
    &= \frac{1}{4} + \frac{1}{16} + \frac{1}{64} + \cdots\\
    &= \frac{1}{4}(1 + \frac{1}{4} + \frac{1}{16} + \cdots)\\
    &= \frac{1}{4} \cdot \frac{1}{1-\frac{1}{4}}\\
    &= \frac{1}{4} \cdot \frac{4}{3}\\
    &= \frac{1}{3}
    \end{align*}
\end{itemize}

\textbf{Step 5: Verify}
\begin{itemize}
    \item Check: Geometric series with $a = \frac{1}{4}$, $r = \frac{1}{4}$
    \item Sum: $\frac{a}{1-r} = \frac{1/4}{3/4} = \frac{1}{3}$ \checkmark
    \item Answer: $\boxed{\textbf{(D)}\ \frac{1}{3}}$
\end{itemize}

\subsection*{Alternative Approach: Geometric Probability}
\textbf{Visualization Method:}
\begin{itemize}
    \item Draw unit square $[0,1] \times [0,1]$ representing $(x,y)$ pairs
    \item Favorable region: union of squares where floor values match
    \item Each square: $[\frac{1}{2^{k+1}}, \frac{1}{2^k}) \times [\frac{1}{2^{k+1}}, \frac{1}{2^k})$ with area $\frac{1}{4^k}$
    \item Total favorable area: $\sum_{k=1}^{\infty} \frac{1}{4^k} = \frac{1}{3}$
    \item Probability: $\frac{\text{favorable area}}{\text{total area}} = \frac{1/3}{1} = \frac{1}{3}$
\end{itemize}

\subsection*{Comprehensive Pitfall Guide}
\begin{enumerate}
    \item \textbf{Sign Error}: Remember $\log_2 x < 0$ for $x \in (0,1)$. The floor values are negative: $-1, -2, -3, \ldots$
    
    \item \textbf{Interval Boundaries}: Use $[2^{-k}, 2^{-k+1})$ (left-closed, right-open), not $(2^{-k}, 2^{-k+1}]$. The floor function makes the left endpoint inclusive.
    
    \item \textbf{Series Indexing}: Our series starts from $k=1$, not $k=0$, because the first interval is $[\frac{1}{2}, 1)$ (corresponding to $k=1$)
    
    \item \textbf{Independence}: Since $x$ and $y$ are independent, we must square probabilities: $P(\text{both in interval}) = P(x \text{ in interval}) \times P(y \text{ in interval}) = \ell^2$ where $\ell$ is the interval length
    
    \item \textbf{Geometric Series Formula}: For $\sum_{k=1}^{\infty} \frac{1}{4^k}$, apply the formula for series starting at $k=1$:
    $\sum_{k=1}^{\infty} \frac{1}{4^k} = \frac{1/4}{1-1/4} = \frac{1}{3}$
    Don't use $\frac{1}{1-1/4} = \frac{4}{3}$ (that's for starting at $k=0$)
    
    \item \textbf{Interval Length Calculation}: The length of $[\frac{1}{2^k}, \frac{1}{2^{k-1}})$ is $\frac{1}{2^{k-1}} - \frac{1}{2^k} = \frac{1}{2^k}$, not $\frac{1}{2^{k-1}}$
\end{enumerate}
\end{solutionbox}

\begin{competitionbox}
\subsection*{A. Time Management}
\begin{itemize}[leftmargin=*]
    \item \textbf{Estimated Time}: 3-4 minutes for late-band problem
    \item \textbf{Breakdown}:
    \begin{itemize}
        \item Understanding floor and logarithm behavior: 45 seconds
        \item Identifying interval structure: 60 seconds
        \item Computing first few probabilities: 45 seconds
        \item Recognizing geometric series and summing: 45 seconds
        \item Verification: 15 seconds
    \end{itemize}
    \item \textbf{If stuck}: Move to next problem after 4 minutes; return if time remains
\end{itemize}

\subsection*{B. Problem Prioritization}
\begin{itemize}[leftmargin=*]
    \item This is a probability problem with logarithms -- prioritize if you're strong in both areas
    \item Requires comfort with floor functions and geometric series
    \item If you don't immediately see the interval partition, this problem can be time-consuming
    \item Skip if you're uncomfortable with logarithms or series; return later
\end{itemize}

\subsection*{C. Risk Assessment}
\begin{itemize}[leftmargin=*]
    \item \textbf{Risk Level}: Medium-High
    \item \textbf{Confidence needed}: Must understand floor function behavior and geometric series
    \item \textbf{Abandonment criterion}: If you don't see interval partition pattern within 90 seconds, move on
\end{itemize}

\subsection*{D. Strategic Notes}
\begin{itemize}[leftmargin=*]
    \item This problem rewards pattern recognition (dyadic intervals, geometric series)
    \item The compound function (floor of logarithm) is deliberately complex -- don't panic
    \item Testing small cases ($k=1, 2, 3$) quickly reveals the pattern
    \item Once you identify geometric series with $r = \frac{1}{4}$, the answer is immediate
\end{itemize}
\end{competitionbox}

\section*{Summary}

This problem exemplifies late-band AMC12 probability: it combines multiple concepts (uniform distributions, logarithms, floor functions, infinite series) in a non-obvious way. The key insights are:
\begin{enumerate}
    \item The floor function partitions $(0,1)$ into dyadic intervals
    \item Each interval has length $\frac{1}{2^k}$ for $k = 1, 2, 3, \ldots$
    \item Probabilities form a geometric series: $\frac{1}{4}, \frac{1}{16}, \frac{1}{64}, \ldots$
    \item The sum is $\frac{1/4}{1-1/4} = \frac{1}{3}$
\end{enumerate}

With practice recognizing floor-of-logarithm patterns and geometric series, problems like this become accessible in 3-4 minutes.

\end{document}